% Бүлэг 1 — Онол ба технологын судалгаа

\chapter{ОНОЛ БА ТЕХНОЛОГЫН СУДАЛГАА}
\label{Chapter2}
\pagecolor{white}

\noindent\textit{\textbf{Бүлгийн товч агуулга.} Энэ бүлэгт цаг захиалгын систем хөгжүүлэхэд шаардагдах гурван бүлэг технологийн онолын үндэс болон сонгох шалтгааныг авч үзнэ. 1.1 дэд бүлэгт мобайл аппликейшн хөгжүүлэлтийн арга хандлагуудыг харьцуулж, Flutter-ийн давуу талыг нотолно. 1.2 дэд бүлэгт системийн тасралтгүй ажиллагаа, автомат байршуулалтад шаардагдах үүлэн технологиуд (Docker, AWS EC2/ECR, GitHub Actions)-ыг судална. 1.3 дэд бүлэгт серверийн талын Backend (бэкэнд) технологиор Node.js/Express REST API-г сонгох үндэслэлийг тайлбарлана.}

\bigskip

Өмнөговь аймгийн Спорт цогцолборын цаг захиалгын системийг хөгжүүлэхийн тулд технологийн сонголт нь дараах гурван шаардлагыг хангахад чиглэсэн юм. Иргэдэд хаанаас ч нэвтрэх боломж олгох мобайл технологи; аюулгүй байдал болон бизнесийн логикийг тусад нь боловсруулах Backend (бэкэнд буюу серверийн талын) технологи; системийг тасралтгүй ажиллуулах, автоматаар шинэчлэх үүлэн дэд бүтэц. Энэ бүлэгт уг гурван шаардлагт нийцэх технологи бүрийн онолын үндэс, зориулалт болон яагаад тэр технологийг сонгосон шалтгааныг тайлбарлана.

%-------------------------------------------------------------------------------

\section{Мобайл аппликейшн хөгжүүлэлт}

Спорт заал захиалах үйлчилгээг иргэдэд хялбар болгохын тулд мобайл аппликейшн зайлшгүй шаардлагатай юм. Өнөөдөр хүн бүр ухаалаг гар утас ашиглаж байгаа тул апп нь аль ч цагт, аль ч газраас захиалга хийх боломж олгодог.

Мобайл аппликейшн хөгжүүлэхэд гурван үндсэн арга хандлага байдаг. \textbf{1) Native хандлага}-аар Android (Kotlin/Java) болон iOS (Swift)-д тус тусад нь тусдаа хэлээр хоёр апп бичдэг. Энэ нь хамгийн өндөр гүйцэтгэлийг өгдөг боловч хоёр баг, хоёр дахин илүү хугацаа шаарддаг. \textbf{2) Hybrid хандлага}-аар вэб технологид (HTML, CSS, JavaScript) тулгуурлан WebView-д ажиллуулдаг. Хөгжүүлэх нь хялбар боловч гүйцэтгэл сул байдаг. \textbf{3) Cross-platform хандлага}-аар нэг кодын баазаас Android болон iOS платформд нэгэн зэрэг гаргадаг. Гүйцэтгэлийн хувьд Native-д ойр, хугацааны хувьд Hybrid-д ойр гэсэн дунд зэргийн оновчтой шийдэл юм.

Эдгээр гурван хандлагын ялгааг шалгуур үзүүлэлт тус бүрээр харьцуулж, аль нь Спорт цогцолборын цаг захиалгын системд хамгийн тохиромжтой болохыг доорх хүснэгтэд харуулав.

\begin{table}[h]
\centering
\caption{Мобайл аппликейшн хөгжүүлэлтийн арга хандлагуудын харьцуулалт}
\label{tab:mobile_approach}
\begin{tabularx}{\textwidth}{|p{2.5cm}|X|X|X|}
\hline
\textbf{Шалгуур} & \textbf{Native} & \textbf{Hybrid} & \textbf{Cross-platform} \\
\hline
Хэл/хэрэгсэл & Swift, Kotlin & HTML, JS & Flutter, React Native \\
\hline
Гүйцэтгэл & Хамгийн өндөр & Доогуур & Native-д ойр \\
\hline
Нэг кодын бааз & Үгүй & Тийм & Тийм \\
\hline
UI нийцтэй байдал & Ялгаатай & WebView хязгаарлагдмал & Бүрэн нийцтэй \\
\hline
Хөгжүүлэлтийн зардал & Өндөр (×2 баг) & Бага & Дунд \\
\hline
\end{tabularx}\\[4pt]
{\footnotesize Эх: Зохиогчийн боловсруулалт}
\end{table}

Хүснэгт~\ref{tab:mobile_approach}-аас үзэхэд Cross-platform хандлага нь нэг кодоос хоёр платформд хүрдэг тул цаг, зардлын хувьд хамгийн үр дүнтэй, гүйцэтгэлийн хувьд Native-д ойр, UI нийцтэй байдал бүрэн хангагддаг гэсэн гурван давуу талыг хослуулдаг. Энэхүү ажлын хувьд оюутны бие даасан туршилтаар нэг хүн хөгжүүлэх тул хоёр салангид код бичих Native хандлага боломжгүй, харин WebView-ийн гүйцэтгэлийн хязгаарлалттай Hybrid хандлага захиалгын real-time шинэчлэлд тохиромжгүй гэж үзсэн. Иймд Cross-platform хандлагыг сонгож, түүний хамгийн өргөн хэрэглэгддэг фреймворк болох \textbf{Flutter}-ийг ашигласан.

\subsection*{Flutter сонгох үндэслэл}

\keyword{Flutter} нь Google компанийн 2018 онд гаргасан нээлттэй эхийн UI фреймворк юм \cite{flutter_doc}. Flutter нь Dart хэл дээр бичигдсэн бөгөөд Skia/Impeller дүрслэлийн хөдөлгүүр ашиглан бүх UI-г платформаас үл хамааран өөрөө зурдаг. Энэ нь Android болон iOS-д яг ижил харагдах баталгаа олгодгийн зэрэгцээ нэмэлт platform-specific код бичих шаардлагыг арилгадаг. Хөгжүүлэлтийн явцад Hot Reload функц нь кодын өөрчлөлтийг апп-г дахин эхлүүлэхгүйгээр шууд харуулдаг тул хөгжүүлэлтийн хурдыг мэдэгдэхүйц нэмэгдүүлдэг. Flutter 3.x хувилбар нь Material Design 3 стандартыг бүрэн дэмждэг бөгөөд \code{ThemeData} классаар дамжуулан дизайн системийг нэгтгэн удирдах боломж олгодог \cite{material3}.

Flutter-ийг React Native, Xamarin зэрэг бусад Cross-platform фреймворкоос ялгах гол давуу талууд нь: (1) UI-г өөрөө зурдаг тул платформын төхөөрөмж шинэчлэгдэх үед UI эвдэрдэггүй, (2) widget суурилсан бүтэц нь дахин ашиглагдах compoнентуудыг хялбар бий болгодог, (3) Material Design 3 болон Cupertino аль алинд бэлэн widget багц өгдөг тул цаг захиалгын дэлгэц зэрэг түгээмэл UI бүтцийг хурдан угсрах боломжтой. Эдгээр давуу талуудаар Flutter нь Спорт цогцолборын мобайл аппликейшний шаардлагуудыг бүрэн хангадаг.

\subsection*{Dart хэл}

\keyword{Dart} нь Google-ийн боловсруулсан, объект хандлагат, хатуу төрлийн хэл юм \cite{dart_doc}. Dart нь AOT (Ahead-of-Time) болон JIT (Just-in-Time) хоёуланг дэмждэг тул хөгжүүлэлтийн үед Hot Reload-ийн хурд, продакшн дээр AOT-ийн хурдан гүйцэтгэл хоёуланг хангадаг. \code{async/await} синтакс нь Firebase Firestore болон API дуудлагуудыг уншихад хялбар байдлаар бичих боломж олгодог. Null Safety механизм нь null pointer алдааг compile цагт илрүүлдэг.

\subsection*{Firebase үйлчилгээ}

\keyword{Firebase} нь Google компанийн санал болгодог BaaS (Backend as a Service) платформ бөгөөд мобайл апп хөгжүүлэгчдэд бэкэнд дэд бүтцийг өөрөө тохируулахгүйгээр бэлэн үйлчилгээнүүдийг ашиглах боломж олгодог \cite{firebase_doc}. Энэ ажилд Firebase-ийн гурван үндсэн үйлчилгээг ашигласан.

\keyword{Cloud Firestore} нь NoSQL, баримт жишгийн өгөгдлийн сан юм \cite{firestore_doc}. Бодит цагийн шинэчлэлт, офлайн хадгалалт, автомат масштабчлал зэрэг онцлогтой бөгөөд collection-document бүтэц нь JSON форматтай ойролцоо тул Flutter SDK-аар шууд хандахад хялбар. Тухайн захиалагдсан цагийн хуваарийг бодит цагаар харуулах шаардлага байгаа тул real-time listener-ийг ашиглав.

\keyword{Firebase Authentication} нь хэрэглэгчийн нэвтрэлтийг удирдах бэлэн үйлчилгээ юм \cite{firebase_auth}. Email/нууц үгийн нэвтрэлт ашиглахад хэрэглэгч бүрт өвөрмөц UID үүсгэж, нэвтрэлтийн токен удирддаг. Энэ ажилд Node.js бэкэнд болон Nodemailer (Gmail SMTP)-ийн 6 оронтой баталгаажуулах кодтой хосолсон бөгөөд хэрэглэгчийн email хаягийн үнэн зөвийг баталгаажуулдаг.

\keyword{Firebase Cloud Functions} нь серверлэс тооцооллын үйлчилгээ бөгөөд AI чатботын үндсэн endpoint болгон ашигласан. Cloud Functions нь Anthropic-ийн API түлхүүрийг Flutter кодод задруулахгүйгээр аюулгүй хадгалах боломж олгодог. Хэрэглэгч чат илгээх бүрт Cloud Function дуудагдаж, Claude Haiku API-аас хариулт авч буцаана.

\subsection*{Локал мэдэгдэл ба AI чатбот}

Захиалгын сануулгыг хэрэгжүүлэхэд \keyword{flutter\_local\_notifications} багцыг ашигласан \cite{local_notifications_doc}. Энэ багц нь Firebase Cloud Messaging серверийн тусламжгүйгээр хэрэглэгчийн төхөөрөмжид мэдэгдлийг schedule хийх боломжтой бөгөөд интернет холболтгүй тохиолдолд ч ажилладаг. Захиалга хийх бүрт тухайн захиалгаас 1 цагийн өмнөх цагт сануулга schedule хийгддэг.

AI чатботын хувьд хоёр загвар хосолсон нөөцлөлтийн архитектурыг ашигласан. \keyword{Anthropic Claude Haiku 4.5} нь Firebase Cloud Functions-д суурилсан үндсэн горимын загвар юм \cite{anthropic_doc}. Хурдан, хэмнэлттэй хариулт өгдөг бөгөөд API түлхүүр Flutter кодод задрахгүй Cloud Functions-д аюулгүй хадгалагддаг. \keyword{Groq Llama 3.1-8b} нь Node.js бэкэнд дамжуулан ажиллах нөөц горимын загвар юм \cite{groq_doc}. Groq-ийн LPU (Language Processing Unit) дэд бүтэц нь маш хурдан хариулт өгдөг давуу талтай. Хэрэв Claude Haiku удаан буюу алдаа гарвал Flutter апп автоматаар Groq руу шилжих тул AI чатботын найдвартай байдал нэмэгдэнэ.

%-------------------------------------------------------------------------------

\section{Үүлэн технологи}

Мобайл клиентийн ард ажиллах серверийн программ хангамжийг ямар дэд бүтцэд байршуулах, хэрхэн тасралтгүй ажиллуулах вэ гэсэн шийдвэр нь системийн найдвартай байдлыг шууд тодорхойлдог. Локал сервер ашиглах нь тасралтгүй ажиллагаа болон аль ч газраас хандах боломжийг хязгаарладаг тул үүлэн технологид суурилсан байршуулалт зайлшгүй шаардлагатай болсон. Дараах гурван технологи нь хамтдаа бэкэндийн найдвартай, автоматаар шинэчлэгдэх дэд бүтцийг бүрдүүлдэг.

\subsection*{Docker контейнерийн технологи}

\keyword{Docker} нь аппликейшн болон түүний бүх хамаарлыг тусгаарлагдсан контейнерт багцалдаг нээлттэй эхийн контейнерийн платформ юм \cite{docker_doc}. Контейнерийн технологи нь өнөөгийн серверийн талын хөгжүүлэлтэд хамгийн өргөн хэрэглэгддэг технологийн нэг бөгөөд Docker нь түүний жишиг хэрэгжилт болж тогтсон. Docker-ийг сонгох гол шалтгаанууд нь: (1) \textit{нээлттэй эх, үнэ төлбөргүй}, Community Edition нь арилжааны зориулалтаар чөлөөтэй ашиглагдана; (2) \textit{орчны тогтвортой байдал}, хөгжүүлэгчийн машин дээр ажилладаг контейнер яг адилхан AWS EC2 дээр ажилладаг тул ``миний компьютер дээр ажиллаж байсан'' гэсэн алдаа арилна; (3) \textit{уян хатан байршуулалт}, нэг л Docker дүрс олон үүлэн үйлчилгээнд (AWS, GCP, Azure) шилжүүлж болно; (4) \textit{экосистемийн өргөн дэмжлэг}, Docker Hub, AWS ECR, GitHub Packages зэрэг олон register-тэй нэгдмэл ажилладаг. Эдгээр давуу талуудаар Docker нь энэхүү ажлын бэкэндийг контейнержуулах хамгийн тохиромжтой шийдэл болсон.

\code{Dockerfile} нь Node.js 18 Alpine суурь дүрс дээр package.json суулгаж, эх кодыг хуулан, 3000 портоор сервер эхлүүлдэг тохиргооны файл юм. Alpine суурь дүрс сонгосноор эцсийн дүрсний хэмжээ жижиг байж, татах болон байршуулах хугацааг хурдасгадаг.

\subsection*{Amazon EC2 виртуал сервер}

\keyword{Amazon EC2} (Elastic Compute Cloud) нь Amazon Web Services (AWS)-аас 2006 онд гаргасан \textit{IaaS} (Infrastructure as a Service буюу Дэд бүтцийн үйлчилгээ) ангиллын виртуал сервер үйлчилгээний нэр юм \cite{aws_ec2_doc}. EC2 нь хэрэглэгчид өөрсдийн хэрэгцээнд тохирсон үйлдлийн систем (Linux, Windows), CPU, RAM, дискний хэмжээтэй виртуал сервер түрээслэн ашиглах боломжийг олгодог. Тухайн серверийн тохиргоо, security group буюу сүлжээний хандалт, дискний шифрлэлт зэргийг бүрэн хянах боломжтой тул нарийн нэвтрэлтийн дүрэм болон тусгай тохиргоо шаардсан системд нэн тохиромжтой.

Энэ ажилд Node.js/Express бэкэнд нь \textbf{ap-southeast-1} (Сингапур) бүсийн EC2 t3.micro instance дээр ажилладаг. Сингапур бүсийг сонгосон шалтгаан нь Монгол улсаас харьцангуй ойр байрладаг AWS бүс тул API-ийн хариу цагийг бууруулдаг. t3.micro instance нь хөгжүүлэлтийн шатанд хямд (saanally сар бүр $\sim$\$8 USD), хангалттай гүйцэтгэлтэй гэж үзэгдэв.

\subsection*{Amazon ECR контейнерийн агуулах}

\keyword{Amazon ECR} (Elastic Container Registry) нь AWS-аас гаргасан, Docker дүрсийг найдвартай хадгалж, түгээх контейнерийн \textit{registry} (агуулах) үйлчилгээний нэр юм \cite{aws_ecr_doc}. ECR нь Docker Hub-ийн AWS дотроо хэрэгжсэн хувилбар бөгөөд олон давуу талтай. Эхнийд, AWS IAM эрхийн удирдлагатай нэгдмэл ажилладаг тул дүрсний нууцлал баталгаатай. Хоёрдугаарт, EC2 болон бусад AWS үйлчилгээтэй ижил сүлжээнд байх тул дүрс татан буулгах хурд маш өндөр. Гуравдугаарт, дүрсийн хувилбарыг автоматаар хадгалж, шинэчлэлт хийх боломжтой.

Энэ ажлын ажиллагааны зарчим нь: GitHub Actions-аар build хийсэн Docker дүрсийг ECR-д push хийж, EC2 instance тэр дүрсийг татаж ажиллуулна.

\subsection*{GitHub Actions CI/CD автоматжуулалт}

\keyword{GitHub Actions} нь GitHub платформд суурилагдсан, YAML файлаар тохируулгатай \textit{CI/CD} (Continuous Integration / Continuous Deployment буюу Тасралтгүй нэгтгэл / Тасралтгүй байршуулалт) автоматжуулалтын программын нэр юм \cite{github_actions_doc}. CI/CD нь хөгжүүлэгчдийн оруулсан кодын өөрчлөлтийг автоматаар туршиж, build хийж, продакшн серверт байршуулдаг арга юм. GitHub Actions-ийн давуу тал нь GitHub репозиторитой нэгдмэл, нийтийн репозиторид сар бүр 2000 минут үнэгүй ажиллагаатай (private репозиторид мөн хязгаартай үнэгүй квот) тул жижиг болон дунд хэмжээний төслүүдэд маш тохиромжтой.

YAML файлаар тодорхойлсон workflow нь тодорхой нөхцөл (push, pull request, schedule) биелэхэд автоматаар ажилладаг. Энэ ажилд \code{backend/} хавтасны файл өөрчлөгдөх бүрт CI/CD хоолой автоматаар ажиллаж, Docker image build хийж AWS ECR-д push хийж, EC2-д шинэ хувилбарыг 3--4 минутын дотор байршуулдаг. Хөгжүүлэгч зөвхөн \code{git push} хийхэд л бэкэнд автоматаар шинэчлэгддэг тул хүний гараар сервер дээр нэвтрэх шаардлага байхгүй болдог.

%-------------------------------------------------------------------------------

\section{Backend (бэкэнд) технологи}

\noindent\textit{\textbf{Тайлбар.} Backend (англиар \textit{back-end}, орчуулбал ``серверийн тал'') нь хэрэглэгчид шууд харагдахгүй, харин аппликейшний бизнесийн логик, өгөгдлийн сан, аюулгүй байдалтай холбоотой үйлдлүүдийг гүйцэтгэдэг серверийн талын программ хангамжийг хэлнэ. Энэ бүлэгт цаашид ``бэкэнд'' гэсэн галигласан үг ба ``серверийн тал'' гэсэн орчуулгыг сольж хэрэглэнэ.}

\bigskip

Мобайл апп дан ганцаар бүх логикийг агуулбал API түлхүүр болон нэвтрэлтийн нууцлал клиент кодод задарч аюулгүй байдал эмзэг болно. Мөн email илгээх, захиалгын давхардал шалгах зэрэг серверийн талын логикийг Flutter-т хийх боломжгүй. Иймд Node.js/Express-д суурилсан тусдаа бэкэнд сервер шаардлагатай болсон. Бэкэнд нь Flutter апп болон Firebase хоорондын зуучлагчийн үүрэг гүйцэтгэж, email баталгаажуулалт, захиалгын менежмент зэрэг функцийг нэгтгэн хангадаг.

\keyword{Node.js} нь Chrome V8 JavaScript хөдөлгүүр дээр суурилсан серверийн тал JavaScript ажиллуулах орчин юм \cite{nodejs_doc}. Node.js-ийн event-driven, non-blocking I/O загвар нь олон нэгэн зэрэг холболтыг цөөн нөөцөөр боловсруулах чадвартай. npm экосистем нь Nodemailer, node-cron, Groq SDK зэрэг шаардлагатай бүх багцыг хурдан суулгах боломж олгодог. \keyword{Express.js} нь Node.js дээр суурилсан хөнгөн, уян хатан веб фреймворк юм \cite{expressjs_doc}. Middleware систем нь хүсэлт боловсруулах дамжлагыг зохион байгуулах боломж олгодог бөгөөд Router нь URL чиглүүлэлтийг модуль болгон хуваана.

\keyword{REST API} (Representational State Transfer) нь Roy Fielding (2000) \cite{rest_fielding} боловсруулсан веб үйлчилгээний архитектурын хэв маяг юм. HTTP протоколын GET, POST, PUT, DELETE аргуудыг ашиглан нөөцийг удирддаг. Сервер хүсэлт бүрийг бие даасан гэж үздэг (stateless) бөгөөд нөөц бүрийг URL-ээр илэрхийлнэ.

Системд нийт 8 API endpoint байдаг бөгөөд тэдгээр нь email баталгаажуулалт (2), хэрэглэгчийн удирдлага (1), захиалгын CRUD үйлдэл (3), цагийн хуваарь авах (1), AI чатбот (1) гэсэн таван бүлэгт хуваагдана. Тэдгээрийн нэр, HTTP арга, гүйцэтгэх үүргийг доорх хүснэгтэд нэгтгэн харуулав.

\begin{table}[h]
\centering
\caption{Node.js/Express REST API endpoint-уудын жагсаалт}
\label{tab:api_endpoints}
\begin{tabularx}{\textwidth}{|p{1.5cm}|p{4.5cm}|X|}
\hline
\textbf{Арга} & \textbf{Endpoint} & \textbf{Тайлбар} \\
\hline
POST & \code{/api/auth/send-code} & 6 оронтой баталгаажуулах код email-р илгээх \\
\hline
POST & \code{/api/auth/verify-code} & Оруулсан кодыг шалгаж, email баталгаажуулах \\
\hline
POST & \code{/api/users/create} & Шинэ хэрэглэгч Firestore-д бүртгэх \\
\hline
GET  & \code{/api/bookings/:venueId} & Заалны захиалгуудын жагсаалт авах \\
\hline
POST & \code{/api/bookings/create} & Шинэ захиалга үүсгэх \\
\hline
PUT  & \code{/api/bookings/:id/cancel} & Захиалга цуцлах \\
\hline
GET  & \code{/api/schedule/:venueId} & Цагийн хуваарь болон тогтмол захиалга авах \\
\hline
POST & \code{/api/chat} & Groq Llama AI чатбот хариулт авах (нөөц горим) \\
\hline
\end{tabularx}\\[4pt]
{\footnotesize Эх: Зохиогчийн боловсруулалт}
\end{table}

Хүснэгт~\ref{tab:api_endpoints}-аас үзэхэд endpoint бүр RESTful загварын дагуу нөөцийн төрөл (auth, users, bookings, schedule, chat)-ийн хүрээнд бүлэглэгдэж, HTTP арга нь үйлдлийн утгыг тодорхой илэрхийлж байна. Тухайлбал, шинэ захиалга үүсгэхэд POST, цуцлахад PUT, харахад GET аргыг ашигласан.

%-------------------------------------------------------------------------------

\section*{Бүлгийн нэгтгэсэн технологийн стек}
\addcontentsline{toc}{section}{Бүлгийн нэгтгэсэн технологийн стек}

Энэ бүлэгт авч үзсэн бүх технологиудыг ашиглах түвшин (мобайл клиент, үүлэн дэд бүтэц, бэкэнд сервер), гүйцэтгэх үүрэг болон сонгосон шалтгааныг доорх хүснэгтэд нэгтгэн харуулав. Хүснэгт нь 1.1 (мобайл), 1.2 (үүлэн), 1.3 (бэкэнд) гэсэн дэд бүлгүүдийн дагуу гурван бүлэгт хуваагдана.

\begin{table}[h]
\centering
\caption{Ашигласан технологийн стекийн нэгтгэсэн хүснэгт}
\label{tab:tech_stack}
\begin{tabularx}{\textwidth}{|p{2.2cm}|p{3cm}|p{3.5cm}|X|}
\hline
\textbf{Давхарга} & \textbf{Технологи} & \textbf{Үүрэг} & \textbf{Сонгосон шалтгаан} \\
\hline
\multirow{8}{*}{\parbox{2.2cm}{1.1\\Мобайл}} & Flutter 3.x & Мобайл клиент (Android/iOS) & Нэг кодоос хоёр платформ, Material Design 3 \\
\cline{2-4}
& Dart & Flutter хэл & AOT compile, Null Safety, async/await \\
\cline{2-4}
& Firebase Auth & Нэвтрэлтийн систем & Аюулгүй, Google-ийн үүлэн дэмжлэгтэй \\
\cline{2-4}
& Cloud Firestore & NoSQL өгөгдлийн сан & Real-time listener, офлайн дэмжлэг \\
\cline{2-4}
& Cloud Functions & Serverless AI endpoint & API түлхүүрийг аюулгүй хадгалах \\
\cline{2-4}
& Claude Haiku 4.5 & AI чатбот (үндсэн) & Хурдан, хэмнэлттэй LLM \\
\cline{2-4}
& Groq Llama 3.1-8b & AI чатбот (нөөц) & LPU дэд бүтэц, маш хурдан хариулт \\
\cline{2-4}
& \code{flutter\_local\_notifications} & Локал мэдэгдэл & FCM сервергүйгээр мэдэгдэл \\
\hline
\multirow{3}{*}{\parbox{2.2cm}{1.2\\Үүлэн}} & Docker & Контейнержуулалт & Орчны тогтвортой байдал, нээлттэй эх \\
\cline{2-4}
& AWS EC2/ECR & Үүлэн сервер ба registry & Бүрэн хяналт, Docker дэмжлэг, IAM аюулгүй байдал \\
\cline{2-4}
& GitHub Actions & CI/CD хоолой & Git workflow-той нэгдмэл, автомат байршуулалт \\
\hline
\multirow{1}{*}{\parbox{2.2cm}{1.3\\Бэкэнд}} & Node.js/Express & REST API сервер & Event-driven, JSON API бариход тохиромжтой \\
\hline
\end{tabularx}\\[4pt]
{\footnotesize Эх: Зохиогчийн боловсруулалт}
\end{table}

Хүснэгт~\ref{tab:tech_stack}-ыг 1.1, 1.2, 1.3 дэд бүлгүүдтэй харьцуулахад технологи бүр ямар давхаргад харьяалагдах нь тодорхой болсон. Мобайл клиент дээр Flutter, Dart, Firebase болон AI үйлчилгээний 8 технологи; үүлэн дэд бүтэц дээр Docker, AWS EC2/ECR, GitHub Actions; серверийн тал дээр Node.js/Express гэсэн нийт 12 технологийг хослуулан гурван давхаргат архитектур бүрдүүлжээ. Эдгээр технологиудын онолын үндэс, давуу талыг энэ бүлэгт тогтоосон тул дараагийн 2-р бүлэгт системийн загварчлал, хөгжүүлэлтийн ажилд орох боломжтой.
